
\chapter{M002 --- Primitive Surface Construction}

\section{Stage 4 --- Comparison (Phase 2)}

\subsection{Mathematical Proposition}

Transforming a Miller covector from the conventional reciprocal basis to
the primitive reciprocal basis preserves the represented crystallographic
plane.

\subsection{Implementation}

The implementation applies the verified basis transformation to obtain a
primitive Miller covector before any integer-lattice construction.

\subsection{Comparison}

\begin{center}
\begin{tabular}{|p{0.42\linewidth}|p{0.36\linewidth}|}
\hline
Mathematical requirement & Implementation status\\
\hline
Plane preserved & Implemented\\
Orientation preserved & Implemented\\
Coordinate representation changes & Yes\\
Scientific surface unchanged & Yes\\
\hline
\end{tabular}
\end{center}

\subsection{Evidence}

Current evidence:

\begin{itemize}
\item Mathematical specification (Stage 1).
\item Mathematical analysis (Stage 2).
\item Implementation decomposition (Stage 3).
\item Module algorithm description.
\end{itemize}

\subsection{Audit Result}

The implementation is consistent with the mathematical interpretation of
this phase as a representation transformation between reciprocal-space
coordinate systems. No discrepancy has been identified at this stage.

\subsection{Outstanding Questions}

The detailed reciprocal-space mapping and sign conventions will be
verified during the evidence review (Stage 5).
